\documentclass{article}
\usepackage[T1]{fontenc}
\usepackage{lmodern}
\usepackage{graphicx}
\usepackage{amsmath,amssymb}
\usepackage[a4paper,margin=2cm]{geometry}
\usepackage[absolute,overlay]{textpos}
\setlength{\TPHorizModule}{1mm}
\setlength{\TPVertModule}{1mm}
\usepackage[indonesian]{babel}
\usepackage{hyperref}
\setlength{\emergencystretch}{2em}
% unit: Halaman 7

\begin{document}

% === Halaman 7 (python/native) ===

Superincreasing Knapsack

• Superincreasing knapsack adalah persoalan knapsack yang dapat dipecahkan dalam 

orde O(n) (jadi, polinomial). 

• Ini adalah persoalan knapsack yang mudah sehingga tidak disukai untuk dijadikan 

sebagai algoritma kriptografi yang kuat.

• Jika senarai bobot disebut barisan superincreasing, maka kita dapat membentuk 

superincreasing knapsack. 

• Barisan superincreasing adalah suatu barisan di mana setiap nilai  di dalam barisan lebih 

besar daripada jumlah semua nilai  sebelumnya. 

• Contoh:   \{1, 3, 6, 13, 27, 52\}  → barisan superincreasing, 

    \{1, 3, 4, 9, 15, 25\}     →  bukan barisan superincreasing

7

\end{document}
